% Restored and adapted detailed content from the previous evaluation draft.

% Source: paper/chapters/06-evaluation/benchmarking.tex

\section{Performance Benchmarking}\label{detail-sec:evaluation-benchmarking}

A defining characteristic of the Motivo compiler is its ``zero-runtime'' architecture, which shifts the
computational burden entirely to the compilation phase.
To evaluate the viability of this approach, it is essential to analyze the compiler's execution speed, particularly as
the complexity of the musical composition scales.

\subsection{Theoretical Complexity}\label{detail-subsec:theoretical-complexity}

The compilation pipeline is composed of several sequential passes, each with a distinct computational complexity:

\noindent \textbf{Frontend and Semantic Analysis.} The complexity of parsing and type-checking is proportional to the
number of lines of source code, $O(S)$, where $S$ is the size of the abstract syntax tree.
Because these passes do not unroll loops or instantiate patterns multiple times, they remain extremely fast
regardless of the composition's length.

\noindent \textbf{Lowering Engine.} The unrolling process has a time complexity of $O(E)$, where $E$ is the total number
of generated \texttt{NoteEvent}s. A simple \texttt{loop (1000)} statement will execute its body 1,000 times during
this phase, meaning the lowering time is proportional to the density of the final timeline, not the brevity of the
source code.

\noindent \textbf{MIDI Backend.} The backend must sort the absolute events by their tick values before computing
delta-times.
This stable sort introduces a time complexity of $O(E \log E)$.
Serializing the events into the binary file requires a final linear pass, $O(E)$.

Therefore, the overall time complexity of the compilation process is bounded by the sorting phase in the backend:
$O(E \log E)$.

\subsection{Execution Speed and Real-Time Viability}\label{detail-subsec:execution-speed-and-real-time-viability}

Because the compiler is implemented in highly optimized C++23 without reliance on garbage collection or dynamic dispatch
for AST traversal, the constant factors in the time complexity are exceptionally small.
In practice, the compiler generates compositions containing tens of thousands of note events in a fraction of a second.

This high-performance metric is critical for the utility of Motivo in a real-world scenario.
The interactive playground environment (discussed in previous chapters) relies on the compiler's ability to re-compile
the entire source file on every keystroke.
Because the execution time for typical musical pieces (ranging from hundreds to thousands of events) sits well below the
threshold of human perception (typically $<100$ milliseconds), the composer experiences real-time auditory and visual
feedback without the need for an interpreted runtime.

\subsection{Memory Footprint}\label{detail-subsec:memory-footprint}

The value-oriented design of the AST and IR ensures that the compiler's memory footprint remains minimal.
The use of \texttt{std::variant} supports spatial locality in the CPU cache during AST traversal.
The intermediate representation holds only the necessary atomic data (pitch, start beat, duration, velocity) packed
tightly into \texttt{struct}s. Consequently, even a composition unrolled into large event streams within the
configured implementation limits will consume only a few
megabytes of RAM, allowing the compiler to run seamlessly on low-end hardware or within memory-constrained Docker
containers.

% Source: paper/chapters/06-evaluation/expressiveness.tex

\section{Code Expressiveness and Compression Ratio}\label{detail-sec:evaluation-expressiveness}

A major objective of this domain-specific language is to provide a highly expressive interface for music composition,
reducing the manual effort required to construct complex arrangements.
We evaluate this expressiveness by analyzing the ``compression ratio''---the difference between the amount of
code written
by the composer and the number of low-level MIDI events generated by the compiler.

\subsection{Algorithmic Conciseness}\label{detail-subsec:algorithmic-conciseness}

Traditional MIDI sequencers and digital audio workstations require the composer to manually place every note on a grid.
In contrast, Motivo allows composers to define structural rules and relationships.

Consider the following drum pattern from the \texttt{example.motivo} file:
\begin{lstlisting}[language=C++, caption={A generative drum pattern using control flow.},label={detail-lst:lstlisting4}]
track percussion using drums {
    pattern groove() {
        for (let i = 0; i < 16; i = i + 1) {
            play hihat from i;
            if (i % 8 == 0) { play kick from i; }
            if (i % 8 == 4) { play snare from i; }
        }
    }
    loop (3) { play groove(); }
}
\end{lstlisting}

This 10-line block of code encapsulates a complex rhythm.
The \texttt{for} loop defines the underlying hi-hat grid, while the modulo-based \texttt{if} conditions procedurally
generate the kick and snare syncopation.
The \texttt{loop (3)} statement then instantiates this 16-beat pattern three times across the timeline.

\subsection{The Compression Ratio}\label{detail-subsec:the-compression-ratio}

To achieve the same result in a traditional MIDI editor, the composer would have to manually draw 48 hi-hat notes, 6
kick drum notes, and 6 snare drum notes---a total of 60 individual note events, which translates to 120 low-level MIDI
messages (\texttt{Note-On} and \texttt{Note-Off}).

Motivo abstracts these 120 MIDI messages into just a handful of declarative rules.
The resulting ``compression ratio'' is exceedingly high.
As the composition grows---for example, looping the groove 32 times instead of 3---the size of Motivo source
code remains
completely unchanged, while the output scales to hundreds or thousands of events.

This expressiveness validates the core thesis: by shifting from a linear, event-based representation to a structural,
programming-based representation, the cognitive overhead of composition is dramatically reduced.
The composer is freed to focus on high-level musical architecture, delegating the tedious task of temporal event
placement entirely to the compiler's lowering engine.

% Source: paper/chapters/06-evaluation/comparative_analysis.tex

\section{Comparative Analysis}\label{detail-sec:evaluation-comparative}

To fully evaluate the utility and positioning of Motivo, it must be contextualized within the existing
ecosystem of music programming tools.
This section provides a comparative analysis against two distinct paradigms: live-coding environments (exemplified by
Sonic Pi) and low-level scripting libraries (such as Python's \texttt{mido}).

\subsection{Comparison with Sonic Pi}\label{detail-subsec:comparison-with-sonic-pi}

Sonic Pi is a highly successful live-coding language embedded in Ruby, designed primarily for performance and education.
While both Sonic Pi and Motivo allows music to be expressed as text, their underlying design philosophies lead
to vastly different compositional workflows.

\subsubsection{Execution Model: Real-Time vs. Compiled}
Sonic Pi utilizes a real-time, interpreted execution model.
The musician writes scripts that trigger sounds ``on the fly'' through the SuperCollider synthesis engine.
Time is managed explicitly by the programmer using the \texttt{sleep} function, which halts the execution thread.

In contrast, Motivo is fully compiled and ``zero-runtime''.
There is no concept of a thread sleeping; instead, the \texttt{play} and \texttt{rest} commands define a mathematical
relationship on a static timeline.
While Sonic Pi excels in improvisational settings where the code is meant to be modified while it runs, Motivo excels
in \textit{structural composition}, where the composer requires a support that a complex arrangement will resolve into
a consistently quantized MIDI file every time it is compiled.

\subsubsection{Structural Abstraction}
In Sonic Pi, polyphony and parallel execution are achieved by spawning new threads (e.g., using \texttt{in\_thread}).
This can lead to synchronization issues if the composer does not meticulously manage the timing of each thread.
Motivo resolves parallelism through the \texttt{voice} construct.
Because voices are unrolled at compile-time by isolating the internal cursor, the compiler supports that parallel
musical lines are aligned according to the compiler cursor model at their inception without any manual thread
management or timing math required
from the composer.

\subsection{Comparison with Raw MIDI Scripting}\label{detail-subsec:comparison-with-raw-midi-scripting}

For offline composition, many developers use general-purpose programming languages combined with MIDI libraries, such as
Python's \texttt{mido} or C++'s \texttt{RtMidi}.
While these tools offer absolute control over the output, they operate at an exceptionally low level of abstraction.

\subsubsection{Event Placement vs. Intent}
When using a library like \texttt{mido}, the composer must manually construct \texttt{Note-On} and \texttt{Note-Off}
messages.
Crucially, because the MIDI file format relies on delta-times, the composer is forced to manually calculate the tick
difference between every sequential event across all tracks.

Motivo abstracts this entirely.
A composer writing \texttt{play (C4, E4, G4):2;} in Motivo communicates their \textit{musical intent}---playing a
C-Major triad for two beats.
Motivo's lowering engine and backend automatically handle the translation into absolute beats, the flattening into
individual notes, the sorting by absolute tick, and the final variable-length delta-time encoding.

\subsubsection{Boilerplate and Expressiveness}
A simple four-beat loop of a hi-hat playing eighth notes requires a significant amount of boilerplate in Python:
creating a file, creating a track, appending tempo metadata, and writing a loop that interleaves \texttt{Note-On} and
\texttt{Note-Off} messages with explicit delta times.

Motivo achieves the same result with minimal syntax:
\begin{lstlisting}[language=C++,label={detail-lst:lstlisting3}]
tempo 120;
track using drums {
    loop (8) { play hihat :0.5; }
}
\end{lstlisting}

This stark contrast in verbosity highlights Motivo's success as a domain-specific tool.
By embedding musical semantics (like default durations, rest tracking, and automatic track/channel assignment) directly
into the language design, the compiler reduces cognitive friction and allows the user to operate strictly within the
musical domain.
